
%%%%%%%%%%%%%%%%%%%%%%%%%%%%%%%%%%%%%%
%%%%%%%% PRINCIPALES PAQUETES %%%%%%%%
%      No modificar esta area        %
%%%%%%%%%%%%%%%%%%%%%%%%%%%%%%%%%%%%%%
\usepackage{fancyhdr}
\usepackage{graphicx}
%\usepackage[spanish, es-tabla]{babel}
\usepackage[english, spanish, es-tabla]{babel}
\usepackage[utf8]{inputenc}
\usepackage{color}
%\usepackage{hyperref}
\usepackage[hidelinks]{hyperref} 
\usepackage{wrapfig}
\usepackage{array}
\usepackage{multirow}
\usepackage{adjustbox}
\usepackage{nccmath}
\usepackage{subfigure}
\usepackage{amsfonts,latexsym} % para tener disponibilidad de diversos simbolos
\usepackage{enumerate}
\usepackage{booktabs}
\usepackage{float}
\usepackage{threeparttable}
\usepackage{array,colortbl}
\usepackage{ifpdf}
\usepackage{rotating}
%\usepackage{cite}
\usepackage{stfloats}
\usepackage{url}
\usepackage{listings}
%%%%%%%%%%%%%%%%%%%%%%%%%%%%%%%%%%%%%%%%%%%
%%%%%%%%%%%%%%%%%%%%%%%%%%%%%%%%%%%%%%%%%%%
%%% CREAR Y REESCRIBIR ALGUNOS COMANDOS %%%
%              No modificar esta area     %
%%%%%%%%%%%%%%%%%%%%%%%%%%%%%%%%%%%%%%%%%%%
\renewcommand\IEEEkeywordsname{Palabras clave}
%%%%%%%%%%%%%%%%%%%%%%%%%%%%%%%%%%%%%%%%%%%
%%%%%%%%%%%%%%%%%%%%%%%%%%%%%%%%%%%%%%%%%%%
%%%           CONFIGURACION             %%%
%              No modificar esta area     %
%%%%%%%%%%%%%%%%%%%%%%%%%%%%%%%%%%%%%%%%%%%
% correct bad hyphenation here
\hyphenation{op-tical net-works semi-conduc-tor} %% Con este comando se especifican como pueden seprarse las sílabas adecuadamente en caso una palabra quede en dos lineas diferentes de texto
\graphicspath{ {figs/} }  %%Ruta donde se encuentran las imágenes, que esté vacio indica que las imagenes están dentro de la misma carpeta que contiene el archivo .tex

%%%%%%%%%%%%%%%%%%%%%%%%%%%%%%%%%%%%%%%%%%
%%%   TITULO Y LOGOS DE PRIMERA PAGINA %%%
%%%%%%%%%%%%%%%%%%%%%%%%%%%%%%%%%%%%%%%%%%

% ---- Archivos de los escudos ----
\newcommand{\LogoIzq}{figs/escudos/escudo-unam.png}
\newcommand{\LogoDer}{figs/escudos/escudo-ciencias.png}

% ---- Altura de los escudos ----
% Se pueden ajustar independientemente debido a que las imágenes
% pueden tener diferentes proporciones o márgenes internos.
\newcommand{\AlturaLogoIzq}{3cm}
\newcommand{\AlturaLogoDer}{3.2cm}


% ---- Construcción del título de la primera página ----
%
% Distribución horizontal:
%
% | 18 % logo UNAM | 64 % título | 18 % logo Ciencias |
%
% De esta forma los logos permanecen junto a los márgenes mientras
% que el título dispone siempre de una zona independiente.
%
\newcommand{\firstpagetitle}[1]{%
	\makebox[\textwidth][c]{%
		%
		% Logo izquierdo
		\parbox[c]{0.18\textwidth}{%
			\raggedright
			\includegraphics[
			width=0.95\linewidth,
			height=\AlturaLogoIzq,
			keepaspectratio
			]{\LogoIzq}%
		}%
		%
		% Título
		\parbox[c]{0.64\textwidth}{%
			\centering
			\textcolor{black}{#1}\par
		}%
		%
		% Logo derecho
		\parbox[c]{0.18\textwidth}{%
			\raggedleft
			\includegraphics[
			width=0.95\linewidth,
			height=\AlturaLogoDer,
			keepaspectratio
			]{\LogoDer}%
		}%
	}%
}


% ---- Modificar \title ----
%
% En main.tex solamente será necesario escribir:
%
% \title{Nombre de la práctica}
%
\let\OriginalReportTitle\title

\renewcommand{\title}[1]{%
	\OriginalReportTitle{%
		\firstpagetitle{#1}%
	}%
}


% ---- Evitar que Hyperref intente meter los logos en los metadatos PDF ----
\pdfstringdefDisableCommands{%
	\def\firstpagetitle#1{#1}%
}


% ---- Estilo de la primera página ----
% Se mantiene para que no aparezcan encabezados ni pies de página,
% pero los logos YA NO pertenecen a fancyhdr.
\fancypagestyle{firstpage}{%
	\fancyhf{}%
	\renewcommand{\headrulewidth}{0pt}%
}